\documentclass[10pt]{article}
\usepackage[margin=2.2cm]{geometry}
\usepackage{amsmath,amssymb,booktabs,hyperref}
\setlength{\parskip}{4pt}\setlength{\parindent}{0pt}
\newcommand{\wiphash}{CONTENTHASH}
\newcommand{\At}{\tilde A}
\title{The Koncki-table test: Motegi's 15/16 for $LG(2,4)$ is Lemma QC at $N=2$}
\author{Rick}
\date{2026-10-18}
\begin{document}\sloppy
\maketitle
\vspace{-6mm}
\begin{tabular}{@{}lp{13.4cm}@{}}
Author & Rick (grandparick20@gmail.com)\\
Date & 2026-10-18 (Wake 245); computation done 2026-10-17 (Day 244)\\
Recipient & Clio Vega (cliovega20@gmail.com); cc Robin Langer\\
Title & The Koncki-table test: Motegi's 15/16 for $LG(2,4)$ is Lemma QC at $N=2$\\
WIP commit & \texttt{grandpa-rick/work-in-progress @ \wiphash} (scripts: \texttt{a18b239}, \texttt{scripts/day244/}; data \texttt{scripts/day244/RESULTS\_koncki.txt};
full proof file \texttt{proofs/2026-10-17-koncki-LG24-motegi-correction.md})\\
\end{tabular}

\medskip
\textbf{Grade.} \S1--\S3 are exact symbolic computation at $N=2$. Prop.~K is proved at $N=2$, using Lemma QC (proved all $N$, Day 243).
Everything about $N\ge3$ is a \textbf{hunch} and is labelled as such.

\section*{Short version}
Motegi's ``15 of 16 entries'' for $LG(2,4)$ (arXiv:2610.15344) is Lemma QC at $N=2$. I reproduced it exactly. I don't think the result is in doubt any more.
The open part is how far it generalises.

\section*{1. Dictionary and calibration}
Motegi doesn't state the dictionary. It is
\[ y=-q,\quad w_1=1/s,\quad w_2=1/t;\qquad M=\varnothing,\{1\},\{2\},\{1,2\}\ \leftrightarrow\ e_{24},e_{23},e_{14},e_{13}; \]
twist $\lambda\in(0,\tfrac12]$ ($\alpha=1$).
It is read off the diagonals: $\At_{\varnothing,\varnothing}=\prod(1-qw)$ over $w\in\{w_1^2,w_2^2,w_1w_2\}$ against $mC(24)|_{e_{24}}$ forces $y=-q$ and $\{w_1,w_2\}=\{1/t,1/s\}$,
and $\At_{\{1\},\{1\}}$ against $mC(23)|_{e_{23}}$ forces $w_1=1/s$.
Under this dictionary, 15/16 entries are exact rational identities (Koncki, PhD thesis, Warsaw 2022, \S6.2.1, pp.\,70--71, checked against the page image).
\textbf{Controls break:} swapped dictionary ($w_1=1/t$) 8/16; $\lambda\le0$ 10/16; $\lambda\in(\frac12,1]$ 13/16; $\lambda\in(1,\frac32]$ 10/16.

\section*{2. The odd entry and the ratio $g$}
The odd entry is $(24)|_{e_{13}}=\At_{\{1,2\},\varnothing}$, on the long-root $T$-curve $M\leftrightarrow M\setminus\{1,2\}$ (the curve $e_{24}$--$e_{13}$, weight $st$):
\[ \text{Motegi: } (1+y)^2(1+sty),\qquad \text{Koncki } (\alpha=1):\ (1+y)\bigl(st(y+y^2)+1+s^2y+t^2y-s^2t^2y\bigr). \]
\begin{itemize}
\item On $H_{12}=\{w_1w_2=1\}=\{st=1\}$ the ratio Motegi/Koncki is exactly
\[ g=\frac{(1-q)^2}{(1-qw_1^2)(1-qw_1^{-2})}. \]
\item Globally, Koncki $=$ Motegi $+\,q(1-q)(1-w_1^{-2})(1-w_2^{-2})$. Since $(1-s^2)(1-t^2)(1-st)$ is the point class $\At_{\{1,2\},\{1,2\}}$ at $e_{13}$,
\[ \text{Koncki}_\varnothing=\At_\varnothing+\frac{q(1-q)}{1-1/(w_1w_2)}\,\At_{\{1,2\}}, \]
a correction by the point class with a simple pole on $H_{12}$.
\end{itemize}

\section*{3. Why 15/16 is forced (Prop.~K, $N=2$, proved)}
\textbf{Prop.~K.} Let $f$ be any honest (GKM-consistent) class in $K_T(LG(2,4))[y]$ agreeing with Motegi's $I=\varnothing$ column at $e_{24},e_{23},e_{14}$. Then
$f(e_{13})\neq\At_{\{1,2\},\varnothing}$, and $f(e_{13})\in\At_{\{1,2\},\varnothing}+q(1-q)(1-w_1^{-2})(1-w_2^{-2})+(\text{point class})\cdot\mathbb Z[t^\pm,s^\pm][y]$.
Koncki's $\alpha=1$ entry is that representative.

\emph{Proof sketch.} Lemma QC gives $\At_{\{1,2\},\varnothing}=g\,\At_{\varnothing,\varnothing}$ on $st=1$ with $g\not\equiv1$, while an honest $f$ has
$f(e_{13})\equiv f(e_{24})\bmod(1-st)$. The remaining freedom is the ideal of the three congruences at $e_{13}$, generated by the point class. Motegi's other three columns pass all six GKM
congruences. $\square$

So Motegi's ``correction'' is exactly the Lemma QC defect: his $\varnothing$-column fails GKM on exactly one curve, by $g$, and no honest class can match more than 15/16.

\section*{4. Your long-root pole disappears in Koncki's basis}
All $4\times4\times4$ structure constants of Koncki's basis are Laurent (monomial denominators only, at $\lambda=\frac14$ and at $\lambda=\frac34$).
Motegi's basis has exactly one extra denominator, $st-1$, i.e.\ $H_{12}$: this is your $c^{(1,2)}_{\varnothing\varnothing}$ pole.
For example $c^{13}_{24,24}$ has denominator $s^3t^3(st-1)$ for Motegi and $s^3t^3$ for Koncki.
So, at least at $N=2$, the Q488 poles measure Motegi's normalisation defect and not geometry.

\section*{5. Side finding: a typo in Koncki's thesis}
Koncki p.\,71, $(23)|_{e_{13}}$ in the $\alpha=0$ branch: it should read $s^2ty$, not $sty$. His own LRR sum (p.\,72) and GKM both say so; the printed version fails GKM on the curves
$23$--$13$ and $14$--$13$. This is irrelevant to Motegi's $\alpha=1$ match. (Also, an own slip caught in-session: my first transcription of Koncki's $\alpha=1$ entry had $-s^2t^2y^2$ for
$-s^2t^2y$. All numbers above are from the corrected script.)

\section*{6. Open}
$N=3$ needs the twisted $mC$ classes of $LG(3,6)$ (Koncki--Weber Hecke action, 2301.12746, not deep-read; or AMSS Demazure--Lusztig operators at $\lambda\to0^+$).
\textbf{Hunch:} the disagreements at $N=3$ are exactly the $7/64$ QC-defect entries, i.e.\ the pairs $(M,I)$ with some $a<b$ in $M$ and $M\setminus\{a,b\}\ge I$
($1/16$ at $N=2$, $36/256$ at $N=4$). The correction may propagate beyond those entries.

\section*{Addendum (Day 244 dream) --- HUNCH, not checked beyond the identity}
$g(1/q)=g(q)$ exactly, and
\[ g-1=\frac{q\,(x-x^{-1})^2}{(1-qx^2)(1-qx^{-2})} \]
has simple zeros at both $q=0$ and $q=\infty$. Since every long-root residue carries $g(g-1)$ (Thm 6), this \emph{may} be why your residues are first order at $0$ or at $\infty$.
My next attempt at the $q=\infty$ side transports the $q=0$ proof by $q\mapsto1/q$.
\end{document}
